\documentclass[10pt,a4paper]{amsart}
\usepackage[margin=19mm]{geometry}
\usepackage{amsmath,amssymb,mathtools}
\usepackage[scr=rsfs,cal=euler]{mathalfa}
\usepackage{xfrac}
\usepackage[unicode,hidelinks,bookmarks=false]{hyperref}
\newcommand{\ie}{i.e.\ }
\newcommand{\cf}{cf.\ }
\newcommand{\resp}{respectively\ }
\newcommand{\Subsection}{\subsection}
\providecommand{\nobreakdash}{\nobreak\hbox{-}}
\setlength{\emergencystretch}{2em}
\multlinegap0pt
\usepackage{amssymb}
\usepackage{mathtools}
\usepackage[shortlabels]{enumitem}
\newtheorem{lemma}{Lemma}
\newcommand{\CAlgfd}{\mathbf{C}^{*}\mathbf{Alg}_{\mathrm{fd}}}
\DeclareMathOperator{\Center}{Z}
\DeclareMathOperator{\Hom}{Hom}
\newcommand{\Htopos}{\mathbf{H}_{\mathbb{Q}}}
\DeclareMathOperator{\Qpot}{Q^{\diamond}}
\title{\Large\bfseries The Degeneracy of the Centre Comonad Model\\ and the Precomposition Obstruction\\ for Quantum Modalities on Presheaf Topoi}
\author{Joey Woo}
\date{Source: arXiv:2606.09467v1 (CC BY 4.0)}
\begin{document}
\maketitle

\noindent\emph{Source and licence.} \Large\bfseries The Degeneracy of the Centre Comonad Model\\ and the Precomposition Obstruction\\ for Quantum Modalities on Presheaf Topoi. Version arXiv:2606.09467v1 (CC BY 4.0), distributed by arXiv under the Creative Commons Attribution 4.0 International licence, http://creativecommons.org/licenses/by/4.0/, as recorded in the arXiv licence metadata for this exact version and shown on the abstract page https://arxiv.org/abs/2606.09467v1. Full licence notice: \texttt{licenses/arxiv-2606.09467-CC-BY-4.0.txt}.\par\smallskip

\noindent\textit{Editorial scope.} Counted unit: the article's lemma lem:annihilation (source0.tex lines 122--128). The article's complete original proof of it is delivered with it (source0.tex lines 130--136, one \texttt{proof} environment of 7 lines). No mathematics is written, repaired or re-derived: the statement and the proof are the article's own lines, in the article's own notation, and no retained line is substituted or re-flowed.  No further source-local context is needed: every cross-reference of every retained line resolves inside the two delivered spans.  Nothing was omitted from any retained span, so the package carries no omission record.  No retained line contains a citation, so the wrapper carries no bibliography and no bibliography entry.  No retained line contains a figure environment, an image-inclusion command or drawing code, so no image is delivered and the package declares no asset dependency.  Editorial formatting only, in addition: an AMS \texttt{amsart} preamble that restates the article's own declarations and macros used by the retained lines, namely the environments \texttt{lemma} on separate counters, together with the standard packages those lines need.  The retained environments print in this document as Lemma 1 in order of appearance, whereas the article prints the same items with the numbering of its own full document.  The complete native source of the article as distributed in the arXiv e-print for this exact version is bundled unchanged as \texttt{sources/202404/source0.tex}.
\begin{lemma}[Annihilation]\label{lem:annihilation}
Let \(A\in\CAlgfd\) be a simple non‑commutative algebra (i.e., \(A \cong M_n(\mathbb{C})\) with \(n>1\)).  Then for every \(B\in\CAlgfd\),
\[
\Qpot(\mathbf{y}_A)(B) = \varnothing,
\]
i.e., \(\Qpot(\mathbf{y}_A)\) is the empty presheaf (the initial object of \(\Htopos\)).
\end{lemma}

\begin{proof}
By definition,
\[
\Qpot(\mathbf{y}_A)(B) = \mathbf{y}_A(\Center(B)) = \Hom_{\CAlgfd}(A,\Center(B)).
\]
The algebra \(\Center(B)\) is commutative.  Suppose there exists a unital \(*\)-homomorphism \(\phi : A \to \Center(B)\).  Since \(A\) is simple, \(\ker\phi\) is a closed two‑sided ideal; by simplicity, it is either \(\{0\}\) or \(A\).  If \(\ker\phi=\{0\}\), then \(\phi\) is injective, implying that \(A\) is a subalgebra of a commutative algebra, hence commutative, contradiction.  If \(\ker\phi = A\), then \(\phi\) is the zero map, which is not unital.  Therefore no unital \(*\)-homomorphism can exist.  Thus the hom‑set is empty for all \(B\). \qedhere
\end{proof}

\end{document}
